% ============================================
% A. Coalescence Outcome Classification
% ============================================
\subsection*{A. Coalescence Outcome Classification}
\label{sec:suppl:classification}

We represent each coalescence outcome in the parental similarity plane with coordinates:
\begin{equation}
s_A = \text{Sim}(C, A), \quad s_B = \text{Sim}(C, B).
\end{equation}

To classify outcomes, we convert these Cartesian coordinates into polar-like variables. We define Radius ($r$) and Selection Preference Ratio (SPR) as follows:

\textbf{Radius ($r$):}
\begin{equation}
r = \sqrt{s_A^2 + s_B^2}
\end{equation}
which measures the overall similarity of the coalesced community to both parents combined. High $r$ indicates strong similarity to at least one parent, while low $r$ indicates poor resemblance to either parent.

\textbf{Selection Preference Ratio (SPR):}
\begin{equation}
\text{SPR} = \frac{\left| \arctan(s_B / s_A) - \frac{\pi}{4} \right|}{\pi/4}
\end{equation}
normalized between 0 and 1. This measures the selection preference of the coalesced community toward one parent versus the other: $\text{SPR} = 0$ indicates balanced contributions from both parents, while $\text{SPR} = 1$ indicates complete preference for one parent.

\textbf{Dominance (D):}
\begin{itemize}
    \item Condition: High $r$ (above threshold $r_{\text{dom}}$)
    \item Strong selection preference: $\theta \leq \theta_{\text{low}}$ or $\theta \geq \theta_{\text{high}}$
    \item Interpretation: The coalesced community is very similar to one parent but not the other.
\end{itemize}

\textbf{Mixture (M):}
\begin{itemize}
    \item Condition: High $r$ (above $r_{\text{mix}}$)
    \item Symmetric angle: $\theta \approx 0.25$ (within tolerance $\Delta\theta$)
    \item Interpretation: The coalesced community maintains comparable contributions from both parents.
\end{itemize}

\textbf{Restructuring (R):}
\begin{itemize}
    \item Condition: Low $r$ (below threshold $r_{\text{res}}$)
    \item Interpretation: The coalesced community is dissimilar to both parents, forming a novel composition.
\end{itemize}
